\begin{textsample}{POS Dim 9 – global\_north\_2024\_04 – Score 2.00 – global\_north\_2024\_04/200376089.txt}  \label{ex:f9_pos_179}
Democratic Sen.
Van Hollen said it still wasn't clear where President Biden stood in continuing to provide aid to Israel, after the president warned Netanyahu that U.S. policy could change if the humanitarian situation in Gaza didn't improve.
Sen.
Chris Van Hollen, D-Md., revealed on Sunday that he still isn't fully clear on the White House ' s position concerning aid and Israeli accountability amid the humanitarian situation in Gaza.
CBS ' Margaret Brennan asked Van Hollen if he was clear on Biden ' s position of conditioning aid to Israel, and noted the Democratic senator has been calling on the White House to"act on the president ' s own standards"regarding national security."I ' m not clear,"he responded."I was glad to see the president, at least as was reported out, finally say to President Netanyahu that if you don't follow these -- my requests -- that there will be consequences, but the president and the White House have yet to lay out what consequences they have and what they want to impose @ @ @ @ @ @ @ @ @ @, the president has made requests to the Netanyahu government, they have ignored those requests, and we have sent more 2,000-pound bombs.
We can not revert back to that.
We have to make sure that when the president requests something, that we have a means to enforce it,"Van Hollen said.
Chris Van Hollen spoke to Margaret Brennan on Sunday about the White House ' s position on Israel. ( Screenshot/CBS ) Brennan asked Van Hollen about the \$14 billion aid package for Israel held up in the House of Representatives, and questioned whether it was being reconsidered.
Van Hollen explained that the president ' s national security \textbf{memorandum} No. 20 stated specifically that"if a recipient of U.S. military assistance, including the Netanyahu government, is restricting the delivery of humanitarian aid, that we should not be sending more weapons.""So it ' s very important that the Biden administration enforce its own policy,"he continued."That was signed by the president of the United States as @ @ @ @ @ @ @ @ @ @ enforced."Brennan also interviewed White House national security communications adviser John Kirby, who was pressed on why the president hasn't spoken out about potential conditions on aid to Israel.
President Joe Biden speaks at a campaign event at Pullman Yards on March 9, 2024, in Atlanta, Georgia. ( Megan Varner/Getty Images )"I ' m not going to get ahead of the president or decisions he might or might not make going forward,"Kirby told CBS."He was very clear in his call with the prime minister that if we don't see some changes in their policies in Gaza and the way they ' re prosecuting operations, we ' re going to have to make some changes of our own."Biden warned Israeli Prime Minister Benjamin Netanyahu on Thursday that U.S. policy in Gaza could change if the Israeli military doesn't do more to improve the humanitarian situation.
During a phone call with his Israeli counterpart, Biden stressed that Israel ' s strikes on"humanitarian workers and @ @ @ @ @ @ @ @ @ @ a White House readout of the call, Fox News Digital previously reported.
President Joe Biden speaks during a meeting of the National Infrastructure Advisory Council in the \textbf{Indian} Treaty Room of the White House in Washington, DC, US, on Wednesday, Dec. 13, 2023. ( Al Drago/Bloomberg via Getty Images ) The same sentiment was echoed later Thursday by Kirby and Secretary of State Antony Blinken, with Blinken telling reporters that the United States would shift gears"if we don't see the changes that we need to see,"according to the Jerusalem Post.
The 30-minute call came after seven aid workers with the World Central Kitchen were killed by Israeli airstrikes this week, adding to concerns about the humanitarian situation in Gaza.
Fox News Digital ' s Michael Lee contributed to this report.
Hanna Panreck is an associate editor at Fox News.
Get all the stories you need-to-know from the most powerful name in news delivered first thing every morning to your inbox

% matched lemmas: indian, memorandum
\end{textsample}
